\documentclass[11pt]{article}
\usepackage{fontspec}
\setmainfont{Times New Roman}
\setsansfont{Arial}
\setmonofont{Consolas}
\usepackage{amsmath,amssymb,mathtools}
\usepackage{geometry}
\usepackage{microtype}
\geometry{a4paper,margin=2.35cm}
\setlength{\parindent}{1.25em}
\setlength{\parskip}{0pt}
\makeatletter
\renewcommand\footnotesize{\@setfontsize\footnotesize{7.7}{8.7}\selectfont}
\makeatother
\emergencystretch=100em
\allowdisplaybreaks
\sloppy
\hbadness=10000
\vbadness=10000
\hfuzz=2pt
\vfuzz=10pt
\raggedbottom

\begin{document}

% Унит: Папер31 Сецтион05 ентры 1 в001. Дирецт Интерславиц Латин транслатион фром the Герман финал-аудитед соурце слице, соурце линес 307--317 / сегментс P31-С0065--P31-С0068, експандед агаинст принтед сцан/ОЦР витнесс паге 95. Фоотноте цоунтер цонтинуес афтер Папер31 Сецтион04 ентры 8.
\setcounter{footnote}{26}

\subsection*{§5. Директне сумы класов}

\textbf{1. \emph{Директна сума класа идеалов или класа представјень.}} Ако идеал $\mathfrak a$ јест директна сума, $\mathfrak a=\mathfrak b+\mathfrak c$, тогда из изоморфизма следује, же всакы идеал того класа такоже става директна сума,
$\bar{\mathfrak a}=\bar{\mathfrak b}+\bar{\mathfrak c}$. При том $\bar{\mathfrak b}$ јест соответно изоморфно к $\mathfrak b$, а $\bar{\mathfrak c}$ к $\mathfrak c$, ако идеалы разматривајут се како $\mathfrak R$-модулы; зато можно говорити о директној суме класа идеалов и о компонентных класах.

Ако колцо $R_{\mathfrak a}$ некојего класа представјень јест директна сума подколц, кторе једновременно сут идеалы в $R_{\mathfrak a}$, тогда заради изоморфизма то само важи за всако колцо $R_{\mathfrak b}$ того класа представјень, при чем компоненты ставајут еквивалентне колца. Зато можно говорити о директној суме класа представјень и о компонентных класах.

Директној суме класа идеалов одговарја директна сума класа представјень; в том смыслу будемо говорити о директној суме класа.\footnote{Питанье, в каквом обсегу все директне сумы класов представјень могут быти створјене директными сумами класов идеалов, туто знову не разсматрја се.} Нај $\beta_1,\ldots,\beta_m$ буде база $\mathfrak b$, а $\gamma_1,\ldots,\gamma_n$ база $\mathfrak c$; понеже сума јест директна, $\beta$ и $\gamma$ заједно творјат линеарно независну базу $\mathfrak a$. Ако $R_{\mathfrak a}$ јест матрично колцо створјено посредством тој базы, тогда матрица из $R_{\mathfrak a}$, приредјена либовольному елементу $\delta$ из $\mathfrak R$, имаје форму
\[
        \begin{pmatrix}D^{\mathfrak b}&0\\0&D^{\mathfrak c}\end{pmatrix}.
\]
Зато определимо $R^{\mathfrak b}$ како систем всих матриц
\[
        \begin{pmatrix}D^{\mathfrak b}&0\\0&0\end{pmatrix}
\]
и соответно определимо $R^{\mathfrak c}$. Треба показати, же $R^{\mathfrak b}$ и $R^{\mathfrak c}$ творјат подколца в $R_{\mathfrak a}$, и же $R_{\mathfrak a}$ јест јих директна сума; једновременно $R^{\mathfrak b}$ и $R^{\mathfrak c}$ ставајут изоморфне к колцам $R_{\mathfrak b}$ и $R_{\mathfrak c}$, створјеным соответно из $\mathfrak b$ и $\mathfrak c$.

Бо чрез хомоморфизм $R^{\mathfrak b}$ одговарја всим таким елементам $\delta$ из $\mathfrak R$, за кторе $\delta\mathfrak c$ изчезаје, то јест кторе принадлежет к идеалному квоциенту нулового идеала чрез $\mathfrak c$, $(0):\mathfrak c$; заради хомоморфизма $R^{\mathfrak b}$ тым самым јест идеалом в $R_{\mathfrak a}$, и то само важи за $R^{\mathfrak c}$. Колцо $R_{\mathfrak a}$ јест сума тыцх идеалов, и имено директна сума, бо продукт $R^{\mathfrak b}R^{\mathfrak c}$ двух колц изчезаје. Наконец, приредјенье, кторе матрици $D^{\mathfrak b}$ приредјује матрицу из $R^{\mathfrak b}$, содржајучу $D^{\mathfrak b}$ и нулы, даје изоморфизм меджу $R^{\mathfrak b}$ и $R_{\mathfrak b}$; в том смыслу компонентным класам $(\mathfrak b)$ и $(\mathfrak c)$ класа идеалов $\mathfrak a$ одговарјајут компонентне класы колца $R_{\mathfrak a}$.

\end{document}

